\section{The criterion and the rational-function input}\label{sec:inputs}

We first state the criterion that turns linear forms into a dimension
bound, then the rational-function construction from which the forms
are built. Both are due to Fischler. His third input, a multiplicity
estimate, will be stated in Section~\ref{subsec:multiplicity} when the
rank proof calls for it.

\subsection{A refined linear-independence criterion}

The rows of the matrix below are the coefficients of integer linear
forms in the target numbers. The two size bounds measure the growth of
those coefficients and the decay of the forms when evaluated at the
target vector. The rank condition is the third hypothesis to be proved.

\begin{proposition}[{\cite[Proposition~1]{F26}, over $\Q$}]
\label{prop:criterion}
Let $\theta=(\theta_0,\ldots,\theta_v)^T$ be a real vector with
$\theta_0\neq0$.
For each $n$ in an infinite set, let $L_n$ be an integer matrix with
$v+1$ columns and any finite number of rows.
Suppose that, for fixed $B>1$ and $\alpha<1$ with $\alpha>0$,
\[
 \max_{i,j}|(L_n)_{i,j}|\leq B^{n+o(n)},\qquad
 \max_i|(L_n\theta)_i|\leq\alpha^{n+o(n)}.
\]
Suppose also that, for sufficiently large $n$,
\begin{equation}\label{eq:weak-rank}
 L_nx=0,\quad x\in\Q^{v+1}\quad\Longrightarrow\quad x_0=0.
\end{equation}
Then
\[
 \dim_\Q\Span_\Q\{\theta_0,\ldots,\theta_v\}
       \geq1-\frac{\log\alpha}{\log B}.
\]
\end{proposition}

Condition \eqref{eq:weak-rank} says that the first column is not a
rational linear combination of the other columns. Full column rank is
not required. This distinction is useful here because the rational
function in Proposition~\ref{prop:rational} is obtained by an existence
argument, and its coefficients need not give full rank.
In our application, $\theta_0=1$ and the other coordinates are even beta
values. We will bound the coefficients and values of the forms in
Section~\ref{sec:estimates}, then prove \eqref{eq:weak-rank} in
Section~\ref{sec:rank}. No lower bound on the value of each individual
form is required.

\subsection{A rational function with small coefficients}

The construction supplies three properties: integer partial-fraction
coefficients of controlled size, exact vanishing of the first Laurent
coefficients at infinity, and extra smallness of subsequent Laurent
coefficients. The first property controls the coefficients of our forms.
The other two have different roles: the vanishing order, governed by
$\omega$, leaves enough decay to differentiate the sums and supplies
zeros for the rank argument; the additional smallness, governed by
$\Omega$, makes the resulting forms exponentially small.

\begin{proposition}[{\cite[Theorem~4]{F26}}]\label{prop:rational}
Let $a$ be a positive integer and let $\omega,\Omega,r$ be rational
numbers with
\[
 a>\Omega\geq\omega\geq1,\qquad r\geq1.
\]
For every sufficiently large integer $n$ such that $rn,\omega n,\Omega n$
are integers, there are integers $c_{i,j}$, not all zero, for
$1\leq i\leq a$ and $0\leq j\leq n$, with the following properties.
Set
\[
 F_n(t)=\sum_{i=1}^a\sum_{j=0}^n\frac{c_{i,j}}{(t+j)^i}
       =\sum_{\nu\geq1}A_\nu t^{-\nu}.
\]
Then $A_\nu=0$ for $\nu<\omega n$,
\begin{equation}\label{eq:coeff-small}
 |c_{i,j}|\leq \exp(n\log\chi+o(n)),
 \quad
 \log\chi=
 \frac{\omega\log2+3\omega^2+\omega^2\log(a+1)
                +\frac12\Omega^2\log r}{a-\omega},
\end{equation}
and, uniformly in $\nu$,
\begin{equation}\label{eq:laurent-bounds}
 |A_\nu|\leq \exp(n\log\chi+o(n))\,n^\nu\nu^a
 \begin{cases}
 r^{\nu-\Omega n},&\nu<\Omega n,\\
 1,&\nu\geq\Omega n.
 \end{cases}
\end{equation}
\end{proposition}

Here $a$ bounds the pole orders, and $\omega n$ is a lower bound for
the order of vanishing at infinity. The parameter $r$ will determine
the summation cutoff $rn$.
The larger threshold $\Omega n$ gives the stronger estimate in the first
line of \eqref{eq:laurent-bounds}; it does not assert additional exact
vanishing. In the tail estimate it will supply a factor $r^{-\Omega n}$,
separate from the power decay controlled by $\omega n$.
All parameters are fixed as $n$ grows. When we later increase $a$ in
Section~\ref{sec:conclusion}, $\omega$ will remain fixed while $\Omega$
grows on the scale $\sqrt{a\log a}$.

Group the coefficients with the same pole order into polynomials
\[
 P_i(z)=\sum_{j=0}^n c_{i,j}z^j,
 \qquad b=\max\{i:P_i\not\equiv0\}.
\]
The $P_i$ will multiply the special functions obtained by summing $F_n$.
We retain the actual largest index $b$, rather than replacing it by $a$,
because the rank proof needs a nonzero highest coefficient $P_b$.
The integer $b$ may depend on $n$, but always satisfies $1\leq b\leq a$.
